
\section{Estándares de ingeniería}

\subsection{ISO/IEC 27001 (controles como requisitos no funcionales; sin certificación)}

En este trabajo, ISO/IEC 27001 se adopta como marco de referencia para definir controles de seguridad de la información a nivel de diseño e implementación. Su aplicación se limita al contexto académico del prototipo y no implica un proceso de certificación formal.

\subsection{ISO 9001 (gestión de calidad del SDLC y control documental)}

ISO 9001 se utiliza como guía para estructurar prácticas de calidad durante el ciclo de vida del desarrollo de software, con énfasis en trazabilidad de decisiones, consistencia de artefactos y control de cambios.

\subsection{Estándares de documentación, versionamiento y trazabilidad}

Se establecen criterios para documentación técnica, manejo de versiones y registro de evidencias de pruebas, de manera que el proceso de desarrollo y evaluación del prototipo sea reproducible y auditable en el entorno académico.
